\section{Después de Transformer: por qué el enfoque de plataforma es preferible a apostar por una sola arquitectura}

\subsection{El debate central de la entrevista}

El entrevistador preguntó «si se debería construir hardware especializado en torno a Transformer». La respuesta de Jensen es propia del enfoque de plataforma: las arquitecturas futuras seguirán cambiando, e incluso podrían llegar a ser \textit{``barely recognizable as transformers''}. Bajo esta premisa, la especialización excesiva reduce el espacio para la innovación futura.

\begin{importantbox}{La metodología de NVIDIA}
No se trata de apostar a «qué arquitectura ganará para siempre», sino de apostar a «que la comunidad de investigación seguirá inventando nuevas estructuras». Por eso la plataforma debe mantener tres capacidades:
\begin{itemize}
    \item programabilidad (programmability)
    \item escalabilidad (scalability)
    \item compatibilidad con el ecosistema (ecosystem compatibility)
\end{itemize}
\end{importantbox}

\begin{warningbox}{Error común: creer que una plataforma de propósito general siempre es más lenta que un ASIC especializado}
A corto plazo, una solución especializada puede llevar una sola tarea al extremo; pero a largo plazo, la velocidad de iteración de los algoritmos y el costo de migrar entre tareas suelen contrarrestar esa ventaja. En sectores que cambian con rapidez, la flexibilidad de la plataforma suele importar más que un pico de rendimiento aislado.
\end{warningbox}

\subsection{Resumen de la sección}
La esencia del debate sobre Transformer no es «quién es más rápido ahora», sino «quién podrá seguir asumiendo cargas de trabajo desconocidas dentro de diez años». El enfoque de plataforma es más robusto en entornos de alta incertidumbre.

